\HMTDivision{\(K\;\longleftrightarrow\;\Lambda_t\;\longleftrightarrow\;\Lambda_t^*\)}{Géométrie récupérable et réciprocité du couplage}{Déformation positive, dualité réticulaire, structure polaire et récupération orientée de l'action.}
\label{rec:apertura}

La fonction de la constante holographique de couplage arithmético--géométrique se précise en étudiant les transformations que son registre détermine. Sa lecture positionnelle conserve douze composantes ordonnées ; sa réalisation incidentielle distingue une direction ; l’archive bilatérale permet de la reconstruire. La composition de ces opérations conduit à une collection de géométries dont le volume fondamental demeure constant et dont l’inversion produit le réseau dual. Le registre intervient donc aussi bien dans la constitution de la géométrie que dans sa récupération.

Ce développement maintient l’ordre génératif établi dans l’article. Les évaluations additive et multiplicative d’APP, l’orientation du TRIT et les opérations de sélection, de transport et de mémoire du TPK produisent l’état enrichi et sa structure discrète conjointe du continu. Les coordonnées régionales et l’incidence exceptionnelle sont des lectures de cet état. Les réalisations vectorielles, réticulaires et hyperboliques qui suivent reçoivent ces sorties avec leurs systèmes de coordonnées et leurs orientations. La généalogie est conservée lors du changement de représentation.

La distinction entre échelle et forme acquiert ici une expression exacte. Une déformation peut dilater certaines directions et en contracter d’autres tout en conservant le produit total de ses facteurs. Dans l’ellipse d’action, le produit des demi-axes conserve la section d’action. Dans les douze rayons déterminés par \(K\), la compensation conserve le covolume réticulaire. Dans l’opérateur de masse, une décomposition sépare le facteur uniforme de la transformation relative de déterminant un. Chaque conservation possède un objet et une démonstration propres ; les lecteurs permettent de les relier sans confondre leurs domaines.


La réciprocité radiale articule deux constructions complémentaires. Dans le système de coordonnées du double cercle spectral, inverser le rayon équivaut exactement à réfléchir la coordonnée complexe par rapport à la droite de partie réelle \(1/2\). Dans le système de coordonnées modulaire de l’ellipse, échanger les demi-axes inverse le rayon et change le signe de la coordonnée de Klein. Les deux correspondances admettent des formules inverses. Le registre d’orientation distingue la lecture conjuguée d’une modification de la section physique d’action.

La substitution des expressions complètes aux lecteurs angulaires fait apparaître une compensation arithmétique spécifique : le terme linéaire de la section antérieure d’action se compense avec le terme correspondant du contre-angle. La continuation régionale demeure intégrale. Le contraste angulaire, le rayon et la coordonnée hyperbolique sont alors exprimés au moyen des mêmes représentations qui interviennent dans \(\alpha\), avec la constante de couplage au sein de sa clôture dodécaphasique. Réciproquement, la coordonnée radiale orientée retrouve la section de retour et permet de réécrire l’expression électronique complète.

La réalisation spectrale distingue à son tour deux classes d’informations. La dualité des réseaux impose une équation fonctionnelle ; le covolume fixe les pôles et leurs résidus ; la déformation modifie la partie entière de la fonction complétée. L’archive récupère la direction qui détermine cette variation et fournit des bornes explicites pour l’erreur de reconstruction. La conservation et la réversibilité s’expriment ainsi par des opérations qui englobent le registre, la géométrie et leurs lectures analytiques.


\HMTMathematicalPaper
